\documentclass[11pt,a4paper]{article}

\usepackage[utf8]{inputenc}
\usepackage[T1]{fontenc}
\usepackage[spanish,es-noshorthands,es-nodecimaldot]{babel}
\usepackage[margin=2.5cm]{geometry}
\usepackage{graphicx}
\usepackage{float}
\usepackage{booktabs}
\usepackage{amsmath}
\usepackage{enumitem}
\usepackage{caption}
\usepackage[hidelinks]{hyperref}

% Las figuras y tablas se generan en ../output/
\graphicspath{{../output/}}

\setlist[itemize]{itemsep=2pt,topsep=4pt}
\captionsetup{font=small,labelfont=bf}

\begin{document}

\begin{titlepage}
\centering
\vspace*{3cm}

{\LARGE\bfseries Trabajo Práctico N$^{\circ}$5\par}
\vspace{0.4cm}
{\large (Semestre de Otoño 2025)\par}
\vspace{0.3cm}
{\large\bfseries Big Data\par}

\vfill

{\large Profesora: María Noelia Romero\par}
\vspace{0.2cm}
{\large Tutor: Tomás Pacheco\par}
\vspace{0.2cm}
{\large Alumna: Catalina Gallardo\par}

\vfill
\end{titlepage}

% ===============================================================
\section*{Construcción de funciones y sistematización de métodos}
% ===============================================================

El análisis de los modelos de clasificación aplicados a los datos de 2004 y 2024 revela
patrones consistentes y algunas diferencias menores entre ambos años. En términos generales,
los modelos presentan un buen desempeño en exactitud, pero serias limitaciones para
identificar correctamente la clase minoritaria (desocupados), que es la variable de interés.

En 2004, los modelos Logit, LDA, LASSO, AdaBoost y Random Forest no logran predecir
correctamente ningún caso de desocupación. La única excepción parcial es KNN, que identifica
dos casos positivos, aunque a costa de siete falsos positivos. Esta situación se refleja en
las métricas: el \textit{accuracy} (0.95) y el MSE (0.05) se mantienen constantes en casi
todos los modelos, pero tanto la precisión como el \textit{recall} son nulos, salvo en el caso
de KNN. El AUC varía ligeramente: alcanza 0.77 con Random Forest, 0.76 con Logit y LDA, y
desciende a 0.50 con AdaBoost, lo que refleja la falta de capacidad discriminativa de este
último. Todos estos resultados se resumen en el Cuadro~\ref{tab:res2004}.

En 2024, los valores de exactitud mejoran levemente a 0.98 y el MSE disminuye a 0.02, lo que
podría interpretarse como una mejora global en la capacidad predictiva. Sin embargo, esta
mejora se debe a la sobrerrepresentación de la clase negativa, ya que ningún modelo logra
clasificar correctamente a los individuos desocupados. En todos los casos, la matriz de
confusión presenta ceros en la columna de verdaderos positivos, por lo que las métricas de
precisión y \textit{recall} permanecen en cero. No obstante, se observan valores elevados de
AUC en algunos modelos: Logit (0.87), LDA (0.86), LASSO (0.84) y Random Forest (0.81). Esto
indica que los modelos sí asignan probabilidades distintas a las dos clases, pero que bajo el
umbral de decisión estándar (0.5) no se produce una clasificación efectiva. Los resultados se
detallan en el Cuadro~\ref{tab:res2024}.

Las Figuras~\ref{fig:roc2004} y~\ref{fig:roc2024} presentan las curvas ROC correspondientes a
cada año. A pesar de que algunos modelos logran áreas bajo la curva razonables, estos valores
no se traducen en un buen desempeño clasificador, lo cual sugiere que el problema central
reside en el severo desbalance de clases. Este desbalance dificulta el aprendizaje de patrones
útiles para identificar la clase minoritaria.

En conclusión, si bien los modelos muestran un buen rendimiento en métricas globales, todos
fallan en el objetivo principal: detectar a los desocupados. Es posible que se necesite un
modelo más completo y una base de datos con una mayor cantidad de desocupados ---y un $n$
mayor en general--- para lograr una mejor predicción.

% ---------------------------------------------------------------
\begin{table}[H]
  \centering
  \caption{Resultados de los modelos para el año 2004.}
  \label{tab:res2004}
  \resizebox{\textwidth}{!}{%
  \begin{tabular}{llcccccc}
    \toprule
    Modelo & Hiperparámetros & Matriz & Acc. & Prec. & Rec. & MSE & AUC \\
    \midrule
    LogisticRegression         & penalty=None                                                  & [[1381, 0] [66, 0]] & 0.95 & 0.00 & 0.00 & 0.05 & 0.76 \\
    LinearDiscriminantAnalysis & ---                                                           & [[1381, 0] [66, 0]] & 0.95 & 0.00 & 0.00 & 0.05 & 0.76 \\
    KNeighborsClassifier       & n\_neighbors=7                                                & [[1374, 7] [64, 2]] & 0.95 & 0.22 & 0.03 & 0.05 & 0.67 \\
    LogisticRegression (LASSO) & C=0.1, penalty=l1, solver=liblinear                           & [[1381, 0] [66, 0]] & 0.95 & 0.00 & 0.00 & 0.05 & 0.73 \\
    RandomForestClassifier     & max\_features=3, max\_samples=0.25, n\_estimators=25          & [[1378, 3] [66, 0]] & 0.95 & 0.00 & 0.00 & 0.05 & 0.77 \\
    AdaBoostClassifier         & learning\_rate=0.001, n\_estimators=25                        & [[1381, 0] [66, 0]] & 0.95 & 0.00 & 0.00 & 0.05 & 0.50 \\
    \bottomrule
  \end{tabular}}
\end{table}

\begin{table}[H]
  \centering
  \caption{Resultados de los modelos para el año 2024.}
  \label{tab:res2024}
  \resizebox{\textwidth}{!}{%
  \begin{tabular}{llcccccc}
    \toprule
    Modelo & Hiperparámetros & Matriz & Acc. & Prec. & Rec. & MSE & AUC \\
    \midrule
    LogisticRegression         & penalty=None                                                  & [[1328, 0] [25, 0]] & 0.98 & 0.00 & 0.00 & 0.02 & 0.87 \\
    LinearDiscriminantAnalysis & ---                                                           & [[1328, 0] [25, 0]] & 0.98 & 0.00 & 0.00 & 0.02 & 0.86 \\
    KNeighborsClassifier       & ---                                                           & [[1328, 0] [25, 0]] & 0.98 & 0.00 & 0.00 & 0.02 & 0.69 \\
    LogisticRegression (LASSO) & C=0.1, penalty=l1, solver=liblinear                           & [[1328, 0] [25, 0]] & 0.98 & 0.00 & 0.00 & 0.02 & 0.84 \\
    RandomForestClassifier     & max\_features=3, max\_samples=0.25, n\_estimators=75          & [[1328, 0] [25, 0]] & 0.98 & 0.00 & 0.00 & 0.02 & 0.81 \\
    AdaBoostClassifier         & learning\_rate=0.001, n\_estimators=25                        & [[1328, 0] [25, 0]] & 0.98 & 0.00 & 0.00 & 0.02 & 0.50 \\
    \bottomrule
  \end{tabular}}
\end{table}

\noindent\textit{Nota}: de los dos métodos de boosting evaluados (AdaBoost y Gradient
Boosting), la función \texttt{evalua\_multiples\_metodos} reporta únicamente el de menor MSE.
Una celda vacía en la columna de hiperparámetros indica que el modelo se estimó con los
valores por defecto.

% ---------------------------------------------------------------
\begin{figure}[H]
  \centering
  \begin{minipage}{0.32\textwidth}
    \centering
    \includegraphics[width=\textwidth]{{LogisticRegression_04}.png}\\[-2pt]
    {\scriptsize Logit}
  \end{minipage}\hfill
  \begin{minipage}{0.32\textwidth}
    \centering
    \includegraphics[width=\textwidth]{{LinearDiscriminantAnalysis_04}.png}\\[-2pt]
    {\scriptsize LDA}
  \end{minipage}\hfill
  \begin{minipage}{0.32\textwidth}
    \centering
    \includegraphics[width=\textwidth]{{KNeighborsClassifier_04}.png}\\[-2pt]
    {\scriptsize KNN}
  \end{minipage}

  \vspace{0.35cm}

  \begin{minipage}{0.32\textwidth}
    \centering
    \includegraphics[width=\textwidth]{{LogisticRegression(LASSO)_04}.png}\\[-2pt]
    {\scriptsize LASSO}
  \end{minipage}\hfill
  \begin{minipage}{0.32\textwidth}
    \centering
    \includegraphics[width=\textwidth]{{RandomForestClassifier_04}.png}\\[-2pt]
    {\scriptsize Random Forest}
  \end{minipage}\hfill
  \begin{minipage}{0.32\textwidth}
    \centering
    \includegraphics[width=\textwidth]{{AdaBoostClassifier_04}.png}\\[-2pt]
    {\scriptsize AdaBoost}
  \end{minipage}

  \caption{Curvas ROC, 2004.}
  \label{fig:roc2004}
\end{figure}

\begin{figure}[H]
  \centering
  \begin{minipage}{0.32\textwidth}
    \centering
    \includegraphics[width=\textwidth]{{LogisticRegression_24}.png}\\[-2pt]
    {\scriptsize Logit}
  \end{minipage}\hfill
  \begin{minipage}{0.32\textwidth}
    \centering
    \includegraphics[width=\textwidth]{{LinearDiscriminantAnalysis_24}.png}\\[-2pt]
    {\scriptsize LDA}
  \end{minipage}\hfill
  \begin{minipage}{0.32\textwidth}
    \centering
    \includegraphics[width=\textwidth]{{KNeighborsClassifier_24}.png}\\[-2pt]
    {\scriptsize KNN}
  \end{minipage}

  \vspace{0.35cm}

  \begin{minipage}{0.32\textwidth}
    \centering
    \includegraphics[width=\textwidth]{{LogisticRegression(LASSO)_24}.png}\\[-2pt]
    {\scriptsize LASSO}
  \end{minipage}\hfill
  \begin{minipage}{0.32\textwidth}
    \centering
    \includegraphics[width=\textwidth]{{RandomForestClassifier_24}.png}\\[-2pt]
    {\scriptsize Random Forest}
  \end{minipage}\hfill
  \begin{minipage}{0.32\textwidth}
    \centering
    \includegraphics[width=\textwidth]{{AdaBoostClassifier_24}.png}\\[-2pt]
    {\scriptsize AdaBoost}
  \end{minipage}

  \caption{Curvas ROC, 2024.}
  \label{fig:roc2024}
\end{figure}

\end{document}
